\subsection{Compact Selmer localization for the Fermat curve of exponent eleven}
\label{subsec:eleven-compact-localization}

Put
\[
 A=\operatorname{Jac}(y^{11}=x(1-x)),\qquad
 B=\operatorname{Jac}(y^{11}=x(1-x)^2),\qquad
 J=\operatorname{Jac}(X^{11}+Y^{11}+Z^{11}=0),
\]
where smooth projective models are understood.  Both $A$ and $B$
have dimension five.  The rank-one assertion for $B(\mathbb Q)$ is
classical: it is the $(1,2,8)$ entry of the table in
\citet[\S4, p.~219]{GrossRohrlich1978}.  Their introduction also
records $\operatorname{rank}J(\mathbb Q)=6$.  The additional input
needed for compact Selmer localization is control of the full
Tate--Shafarevich group, rather than the rank alone.

\begin{proposition}[The rank-one cyclic factor at eleven]
\label{prop:eleven-rankone-arithmetic}
The abelian variety $B$ has rational Mordell--Weil rank one,
$B(\mathbb Q)_{\mathrm{tors}}\simeq\mathbb Z/11\mathbb Z$, and finite
$\Sha(B/\mathbb Q)$.  In the standard weight-one motivic normalization,
\[
 3.114<L'(B/\mathbb Q,1)<3.137.
\]
\end{proposition}
\begin{proof}
Write $K=\mathbb Q(\zeta_{11})$, $F=K^+$, and
$\lambda=1-\zeta_{11}$.  The Jacobi-sum character for the triple
$(1,2,8)$ is
\[
 j(\mathfrak p)=-\sum_{x\ne0,1}
       \chi_{\mathfrak p}(x)\chi_{\mathfrak p}(1-x)^2,
 \qquad\mathfrak p\nmid11.
\]
Its Hecke $L$-function is $L(B/\mathbb Q,s)$, and every complex
embedding of $j(\mathfrak p)$ has absolute value
$\sqrt{N\mathfrak p}$; see \citet[\S1, pp.~202--203]{GrossRohrlich1978}
and \citet[\S\S2.1--2.2, author pp.~5--6]{ShuExpanded2021}.

The conductor and sign can be checked without numerical inference.
For $r=1,s=2,t=8$ put $u=r^rs^s(t-11)^t$.  Direct arithmetic gives
\[
 u\equiv108\pmod{121},\qquad
 d=\frac{u^{10}-1}{11}\equiv10\pmod{11}.
\]
Thus \citet[Proposition~3.12(2), author p.~16]{ShuExpanded2021}
gives conductor $\lambda^2$.  The root-number formula in
\citet[Theorem~3.13, author pp.~17--18]{ShuExpanded2021} gives
\[
 W=\left(\frac2{11}\right)
       \left[-\left(\frac{-2rstd}{11}\right)\right]=-1,
\]
because $-2rstd\equiv10\pmod{11}$ and both displayed Legendre
symbols are $-1$.  The rational conductor is therefore
$11^9 N\lambda^2=11^{11}$.  On setting
\[
 a=\frac{\sqrt{11^{11}}}{(2\pi)^5},
\]
the entire completed function satisfies
\[
 \Lambda(s)=a^s\Gamma(s)^5L(B/\mathbb Q,s),\qquad
 \Lambda(s)=-\Lambda(2-s).
\]

We give the complete finite nonvanishing calculation.  Let $Y$ be the
product of five independent exponential random variables of mean one,
and define
\[
 H_5(x)=\mathbb E\bigl(\log(Y/x)_+\bigr)
       =\frac1{2\pi i}\int_{(c)}
          \frac{\Gamma(1+z)^5}{z^2}x^{-z}\,dz\quad(c>0).
\]
Here the subscript $+$ means the positive part.  Mellin inversion
follows either by integrating the tail distribution of $Y$ or from
the Mellin transform of the positive logarithm.  If
$L(B/\mathbb Q,s)=\sum a_n n^{-s}$, shifting the contour across
zero and using the odd functional equation gives
\[
 L'(B/\mathbb Q,1)
       =2\sum_{n\ge1}\frac{a_n}{n}H_5(n/a).
\]
Indeed $a^z\Gamma(1+z)^5L(1+z)$ is odd in $z$; the residue after
division by $z^2$ is $L'(1)$, and the two vertical integrals are
equal.  The bounds below justify the absolute convergence and
contour operations.  Since $L(1)=0$, differentiating the completed
function introduces no extra logarithmic term at the center.

Let $b_K(n)$ denote the number of integral ideals of $K$ of norm $n$.
The Hecke expansion gives $|a_n|\le\sqrt n\,b_K(n)$.  For every
$k>0$, the inequality $\log(v)_+\le v^k/k$ gives
\[
 H_5(x)\le\frac{\Gamma(k+1)^5}{k x^k}.
\]
Consequently, if $k+1/2>\sigma>1$, the omitted coefficients beyond
$N$ contribute at most
\begin{equation}
 \frac{2\Gamma(k+1)^5}{k}\,a^k
       N^{\sigma-k-1/2}\zeta_K(\sigma).
 \label{eq:eleven-rankone-tail}
\end{equation}

All coefficients through $N=50000$ are obtained by exact finite-field
counts.  For a prime $p\ne11$ put $f=\operatorname{ord}_{11}(p)$.
For even $f$, \citet[Lemma~1.1, p.~203]{GrossRohrlich1978} gives
\[
 P_p(T)=(1+p^{f/2}T^f)^{10/f}.
\]
For $f=5$, enumerate the logarithms in $\mathbb F_{p^5}^{\times}$
and count the classes of $\log(x)+2\log(1-x)$ modulo eleven.
Only $p=3,5$ are needed at the cutoff; the resulting polynomials are
\[
 P_3(T)=1+9T^5+243T^{10},\qquad
 P_5(T)=1+75T^5+3125T^{10}.
\]
The trace in $K$ is five times the trace over the quadratic field
fixed by $p$-Frobenius, so division by five in this computation is
exact.  At $f=1$, the same enumeration is in $\mathbb F_p$.
If $c_0,\ldots,c_{10}$ are the counts, then
$j=-\sum c_i\zeta_{11}^i$; taking the product over its ten
conjugates, with Newton's identities in integer arithmetic, gives
the required local coefficients.  The bad Euler factor at eleven
is one.  Primes with $p^f>N$ have no nonconstant contribution below
the cutoff.  Euler inversion and multiplicativity give 573 nonzero
coefficients, starting with
\[
 a_1=1,\quad a_{23}=-12,\quad a_{67}=-23,\quad
 a_{89}=-1,\quad a_{243}=-9.
\]

The finite sum is enclosed by the single contour integral
\[
 L'_N=\frac2\pi\Re\int_0^\infty
       \frac{\Gamma(3+it)^5}{(2+it)^2}
       a^{2+it}\sum_{n\le N}a_n n^{-3-it}\,dt.
\]
Put $S_N=\sum_{n\le N}|a_n|n^{-3}$.  From the gamma integral and
its recurrence, $|\Gamma(3+it)|\le\Gamma(3+m)|t|^{-m}$.
The part $t\ge T$ is therefore at most
\begin{equation}
 \frac2\pi a^2S_N
    \frac{\Gamma(3+m)^5}{(5m+1)T^{5m+1}}.
 \label{eq:eleven-rankone-vertical}
\end{equation}
Arb arithmetic at 128 bits encloses the integral on $[0,20]$ by
twenty unit intervals.  The Dedekind-zeta bound is obtained from
\[
 \zeta_K(\sigma)=\zeta(\sigma)
       \prod_{\substack{\chi\bmod11\\\chi\ne1}}L(\sigma,\chi),
 \qquad
 L(\sigma,\chi)=11^{-\sigma}\sum_{a=1}^{10}
                   \chi(a)\zeta(\sigma,a/11),
\]
so that the factor at eleven is included.  The resulting certified
bounds are
\[
 \begin{aligned}
 \zeta_K(11/10)&<3.007268,\\
 |L'_N-3.12547064977|&<3.14\cdot10^{-12},\\
 \text{vertical remainder}&<7.6\cdot10^{-26},\\
 \text{coefficient remainder}&<0.011188.
 \end{aligned}
\]
The last two lines use $T=20,m=17$ and $k=4$, respectively.
These explicit inequalities prove the asserted interval.

We now check the arithmetic theorem's precise hypotheses.
Automorphic induction of $j$ from $K$ to $F$ is a cuspidal
parallel-weight-two Hilbert form $f$: its two infinity exponents
differ by one, so the character is not a norm character.  Its finite
central character is trivial, since the inducing unitary character
restricts to $\eta_{K/F}$ on $\mathbb A_F^\times$, and the determinant
of induction multiplies this by $\eta_{K/F}$ once more
\citep[\S3.4, author p.~15]{ShuExpanded2021}.  One has
\[
 L(f/F,s)=L(B/\mathbb Q,s).
\]
Every coefficient conjugate of $f$ has the same $L$-function, since
conjugating a Jacobi character reindexes the prime ideals of the
abelian field $K$ while preserving their norms.

The CM type is primitive by
\citet[Lemma~1.6, pp.~205--207]{GrossRohrlich1978}: at eleven the
exception involving a nontrivial cube root of unity in
$(\mathbb Z/11\mathbb Z)^\times$ is impossible.  Thus $B$ is
absolutely simple.  The real CM subfield $F$ acts on $B_F$ and is
fixed by Rosati.  Its five coefficient components are the five
two-dimensional inductions of the coefficient conjugates of $j$.
The Hecke field of $f$ is contained in $F$ and is exactly $F$.
For an explicit check, choose $5$ as primitive element of
$\mathbb F_{23}^{\times}$ and send it to $\zeta_{11}$ under the
multiplicative character.  The exact Jacobi-sum coefficient vector is
\[
 (c_0,\ldots,c_{10})=(3,3,2,1,4,2,1,0,0,4,1).
\]
The coefficients of $\zeta_{11}^i+\zeta_{11}^{-i}$, $1\le i\le5$,
in the real trace of the Jacobi sum are therefore
$(4,6,1,4,3)$.  They are not all equal.  Since the sole rational
relation among one and these five cyclotomic pairs is
$1+\sum_{i=1}^5(\zeta_{11}^i+\zeta_{11}^{-i})=0$, this trace is
not rational.  The only proper subfield of $F$ is $\mathbb Q$,
so it generates $F$.  The attached-variety existence assertion
preceding Zhang's Theorem~A and the matching Tate representations
now identify $B_F$ with the abelian variety attached to $f$, up to
isogeny, by \citet[\S5, Corollary~2, p.~22]{Faltings1986}.

The real degree five is odd.  The verified simple zero therefore
meets every hypothesis of \citet[Theorem~A, p.~28; proof in
\S7.2, p.~143]{Zhang2001}, which gives finite full $\Sha(B/F)$.
There is no CM exclusion or additional finite-place hypothesis in
this odd-degree case.  The kernel of
$\Sha(B/\mathbb Q)\to\Sha(B/F)$ lies in
$H^1(\operatorname{Gal}(F/\mathbb Q),B(F))$, a finite group because
$B(F)$ is finitely generated.  Thus $\Sha(B/\mathbb Q)$ is finite.
The rational rank one is the cited classical table entry, and
\citet[Theorem~1.1, p.~202]{GrossRohrlich1978} gives the stated torsion.
\end{proof}

\begin{lemma}[An integral decomposition at eleven]
\label{lem:eleven-integral-isogeny}
There are rational homomorphisms
\[
 \Phi:J\longrightarrow A^3\times B^6,\qquad
 \Psi:A^3\times B^6\longrightarrow J
\]
with $\Phi\Psi=[11]$ and $\Psi\Phi=[11]$.  In particular,
\[
 T_3J\simeq(T_3A)^{\oplus3}\oplus(T_3B)^{\oplus6}
\]
integrally and $G_{\mathbb Q}$-equivariantly.
\end{lemma}
\begin{proof}
For $a=1,\ldots,9$ use the cyclic quotient
$C_a:y^{11}=x(1-x)^a$ and the rational map
\[
 x=-X^{11}/Z^{11},\qquad
 y=(-1)^{a+1}XY^a/Z^{a+1}.
\]
Its geometric deck subgroup is
$H_a=\{(\zeta_{11}^u,\zeta_{11}^v):u+av=0\}$.
The ninety characters in $H^1(X_{11})$ have indices
$(i,j)$ with $i,j,i+j\ne0$ modulo eleven.  They partition
into the nine lines $j=ai$.  The branch-exponent triples
for $a=1,5,9$ are equivalent to $(1,1,9)$ under permutation
and common nonzero scaling modulo eleven; those for
$a=2,3,4,6,7,8$ are equivalent to $(1,2,8)$.
These branch-coordinate and monomial transformations are defined
over $\mathbb Q$, so the product of the nine quotient Jacobians
is $A^3\times B^6$.

For the quotient maps $f_a$, pullback followed by trace satisfies
$f_a^*f_{a,*}=\sum_{h\in H_a}h$, while
$f_{a,*}f_a^*=[11]$.  The first operator acts by eleven
on its own character lines and by zero on every other line.
It follows that $f_{a,*}f_b^*=0$ for $a\ne b$ and
$\sum_a f_a^*f_{a,*}=[11]$.  Faithfulness on a rational Tate
module turns these into identities of abelian-variety
homomorphisms.  Taking the product of the traces and the
sum of the pullbacks gives $\Phi$ and $\Psi$.
Since eleven is a unit in $\mathbb Z_3$, these identities
give the integral Tate-module isomorphism.
\end{proof}

\begin{lemma}[Rank-one factors and integral compact localization]
\label{lem:rank-one-compact-localization}
Let $D/\mathbb Q$ be an abelian variety with finite
$\Sha(D/\mathbb Q)$. Suppose that $D$ is rationally isogenous to
a product $\prod_iD_i$ with
$\operatorname{rank}D_i(\mathbb Q)\le1$ for every $i$.
For every rational prime $\ell$, the integral compact localization
\[
 \operatorname{Sel}(\mathbb Q,T_\ell D)
       \longrightarrow H^1(\mathbb Q_\ell,T_\ell D)
\]
is injective. There is no restriction on the isogeny's degree.
\end{lemma}
\begin{proof}
The inverse-limit Kummer sequence and finite $\Sha$ identify the
source with $D(\mathbb Q)^{\wedge_\ell}$, including its finite
$\ell$-primary torsion. At the real place, $H^1(\mathbb R,D)$ is
killed by two and has zero Tate module, so the integral Tate-module
Kummer condition there is automatic, including when $\ell=2$.
Local Kummer theory embeds
$D(\mathbb Q_\ell)^{\wedge_\ell}$ into the target cohomology group.
It therefore suffices to prove injectivity on these completions.

The logarithm on $D_i(\mathbb Q_\ell)$ is injective on a sufficiently
small open formal subgroup, and a nonzero integer multiple of every
point belongs to that subgroup. Its kernel is consequently exactly
the torsion subgroup. A nontorsion rational point remains nontorsion
locally: an equality $[m]P=0$ over $\mathbb Q_\ell$ between rational
points already holds over $\mathbb Q$. Thus a generator of a rank-one
rational Mordell--Weil group has nonzero local logarithm, proving
injectivity of
\[
 D_i(\mathbb Q)\otimes\mathbb Q_\ell
       \longrightarrow\operatorname{Lie}(D_i)(\mathbb Q_\ell).
\]
For rank zero the source is zero. The product has the same
injectivity, and a rational isogeny induces isomorphisms on the
rationalized Mordell--Weil spaces and on the local Lie algebras.
Hence the asserted localization for $D$ is injective after
inverting $\ell$.

Any remaining integral kernel is therefore contained in
$D(\mathbb Q)[\ell^\infty]$. No nonzero such point disappears in
$D(\mathbb Q_\ell)^{\wedge_\ell}$. Indeed, if
$P\in\bigcap_n\ell^nD(\mathbb Q_\ell)$, choose $Q_n$ with
$\ell^nQ_n=P$. When $P$ is torsion, every $Q_n$ is torsion.
The local torsion subgroup is finite, since it injects into the
finite quotient by a sufficiently small torsion-free formal
subgroup. For $n$ larger than its $\ell$-primary exponent,
$\ell^nQ_n$ cannot be a nonzero $\ell$-primary torsion point.
Thus $P=0$, which eliminates the entire integral kernel.
\end{proof}

\begin{theorem}[Compact localization at exponent eleven]
\label{thm:eleven-compact-selmer}
The full Tate--Shafarevich group $\Sha(J/\mathbb Q)$ is finite,
$\operatorname{rank}J(\mathbb Q)=6$, and for every rational prime
$\ell$ the integral compact Selmer localization
\[
 \operatorname{Sel}(\mathbb Q,T_\ell J)
             \longrightarrow H^1(\mathbb Q_\ell,T_\ell J)
\]
is injective. At $\ell=3$ its source is a free $\mathbb Z_3$-module
of rank six.
\end{theorem}
\begin{proof}
The rank-zero factor $A$ has finite full $\Sha$ and rational
torsion $\mathbb Z/11\mathbb Z$ by Proposition~\ref{prop:eleven}.
The preceding proposition gives the corresponding assertions for
$B$, with rank one.  Finiteness of full $\Sha$ is invariant under
isogeny and finite products, so the integral isogeny lemma gives
finite $\Sha(J/\mathbb Q)$ and rank six.

For a number-field abelian variety $D$, the inverse-limit Kummer
sequence is
\[
 0\longrightarrow D(\mathbb Q)^{\wedge_3}
 \longrightarrow\operatorname{Sel}(\mathbb Q,T_3D)
 \longrightarrow T_3\Sha(D/\mathbb Q)\longrightarrow0.
\]
Finite $\Sha$ makes the last term zero.  Thus the Selmer group of
$A$ is zero and that of $B$ is $\mathbb Z_3$, with no three-primary
rational torsion.  The integral isogeny identifies the Selmer group
of $J$ with six copies of the latter group, respecting all local
Kummer conditions.

Lemma~\ref{lem:rank-one-compact-localization} applies to the
three rank-zero and six rank-one factors. It proves the claimed
integral localization for every $\ell$, including $\ell=11$,
where the isogeny degree need not be a unit.
\end{proof}

\begin{remark}[Compact and affine scope]
\label{rem:eleven-compact-scope}
This proves the compact localization condition for the base exponent
eleven.  It does not establish the corresponding condition for
$11\cdot3^r$ with $r\ge1$, nor the global cuspidality of all affine
pro-three Selmer sections at exponent eleven.  In the latter problem,
the distinct coordinate cusps have differences of order eleven, so
their compact three-adic Kummer classes coincide.  The good-place
argument that selects one cusp using a nonzero three-primary cusp
difference therefore cannot be imported from the exponents divisible
by three.  Separation of the affine cusp packets at three does not,
by itself, impose the needed global boundary support condition.
\end{remark}
